\documentclass[twocolumn, fontsize=10pt]{article}


\raggedbottom
\setlength{\parskip}{0pt}

\usepackage[margin=0.80in]{geometry}
\usepackage{lipsum,mwe,abstract}
\usepackage[T1]{fontenc}
\usepackage[english]{babel}
\usepackage{fancyhdr}
\pagestyle{fancyplain}
\fancyhead{}
\fancyfoot[C]{\thepage}
\setlength\parindent{0pt}
\usepackage{booktabs,tabularx,caption,hyperref}
\usepackage{natbib}
\usepackage{amsmath,amsfonts,amsthm}
\usepackage{graphicx}
\usepackage{float}
\usepackage{subcaption}
\usepackage{comment}
\usepackage{enumitem}
\usepackage{cuted}
\usepackage{booktabs}
\usepackage{sectsty}
\usepackage{dblfloatfix}      % lets figure* go at the bottom [b]
\usepackage{placeins}  % no automatic barriers at \section

\setlength{\dbltextfloatsep}{8pt plus 2pt minus 2pt}
\usepackage{cuted}     % gives the strip environment (full-width block)
\usepackage{capt-of}   % for \captionof outside figure

\usepackage{graphicx}
\usepackage[dvipsnames]{xcolor}
\renewenvironment{abstract}{
 {\small
  \begin{center}
  \bfseries \abstractname\vspace{-.5em}\vspace{0pt}
  \end{center}
  \list{}{\setlength{\leftmargin}{0mm}\setlength{\rightmargin}{\leftmargin}}
  \item\relax}
 }{\endlist}


\usepackage[]{mdframed}
\makeatletter
\renewcommand{\maketitle}{\bgroup\setlength{\parindent}{0pt}
\begin{flushleft}
  \@title
  \@author \\
  \@date
\end{flushleft}\egroup
}

\makeatother

\newcommand{\framedtext}[1]{%
\par%
\noindent\fbox{%
    \parbox{\dimexpr\linewidth-2\fboxsep-2\fboxrule}{#1}%
}%

}

\definecolor{RBlue}{HTML}{2E4A7C}
\definecolor{RBlueLight}{HTML}{90AEC1}
\definecolor{RSand}{HTML}{F2CB80}
\definecolor{RGray}{HTML}{E0E2DC}
\definecolor{CodeBG}{HTML}{F7F8FA}

\definecolor{ElegantNavy}{HTML}{0E2A47}
\definecolor{SteelBlue}{HTML}{4B6E8A}
\definecolor{DeepGrey}{HTML}{333333}


\usepackage[most]{tcolorbox}
\tcbuselibrary{minted}
\usepackage{minted}

\definecolor{RCharcoal}{HTML}{2B2B2B}

\hypersetup{
  colorlinks=true,
  linkcolor=RBlue,
  citecolor=blue,
  urlcolor=RBlue
}

\colorlet{bg}{CodeBG}
\DeclareTCBListing{mintedbox}{O{}m!O{}}{%
  breakable=true,
  listing engine=minted,
  listing only,
  minted language=#2,
  minted style=default,
  minted options={
    linenos,
    breaklines=true,
    fontsize=\small,
    numbersep=8pt,
    #1},
  boxsep=0pt,
  left=25pt,
  top=3pt,
  bottom=3pt,
  arc=0pt,
  colback=white,
  colframe=SteelBlue,
  boxrule=0pt,
  leftrule=3pt,
  enhanced,
  #3
}

\definecolor{captiongray}{gray}{0.35}
\captionsetup[figure]{font=small, labelfont={bf,color=captiongray}, textfont={color=captiongray}}
\captionsetup[table]{font=small,  labelfont={bf,color=captiongray}, textfont={color=captiongray}}
\captionsetup[subfigure]{font=footnotesize, labelfont={bf,color=captiongray}, textfont={color=captiongray}}

\newtcolorbox{myquote}[1][]{%
    colback=black!5,
    colframe=black!5,
    notitle,
    borderline west={2pt}{0pt}{RoyalBlue!80!black},
    enhanced,
    breakable,
    }

\sectionfont{\color{ElegantNavy}\bfseries\fontsize{11}{13}\selectfont\MakeUppercase}
\subsectionfont{\color{SteelBlue}\bfseries\fontsize{10}{12}\selectfont}
\subsubsectionfont{\color{DeepGrey}\bfseries\fontsize{9}{11}\selectfont}

\title{
\Large
\centering
\textbf{Classifying Breast Cancer: A Comparative Analysis} \\
[8pt]
\Large \textcolor{RCharcoal}{Machine Learning, Assignment 2, Supervised Learning Methods} \\[10pt]}
\date{\today}
\author{{\normalsize I\~nigo Arriazu Garc\'ia, Laia Colom\'e Xicoy, Andrea P\'erez Valle, Joana Ros Alonso, Gabriel Torres Zamora }}

\begin{document}
\twocolumn[ \maketitle ]



% ─────────────────────────────────────────────────────────────────────────────
\section{Introduction}
\label{sec:intro}

{Supervised learning} is an ML branch that learns a mapping from input features to their output labels using annotated training examples \citep{hastie2009elements}. Unlike unsupervised approaches, supervised methods can be evaluated directly against known outcomes, making them particularly suited to clinical decision support tasks where ground-truth diagnoses are available.\\

In oncology, it supports early detection by learning to distinguish malignant from benign tissue based on quantitative measurements \citep{bishop2006pattern}. Breast cancer is the most frequently diagnosed cancer in women worldwide and one of the leading causes of cancer mortality \citep{siegel2023cancer}; automated classification tools can contribute to earlier diagnosis and more consistent clinical decision-making.\\

Among supervised models, there is no universally superior classification algorithm: different methods make different assumptions about the shape of the decision boundary, the distribution of the data, and so on \citep{james2021statistical}. 



% ─────────────────────────────────────────────────────────────────────────────
\section{Materials and Methods}
\label{sec:methods}



\subsection{Dataset}

The study uses the {\textit{Breast Cancer Wisconsin (Diagnostic)}} dataset, available directly from \texttt{scikit-learn} via \texttt{load\_breast\_cancer()}. The dataset contains \textbf{569 patient records}, each described by \textbf{30 real-valued features} derived from digitized images of fine needle aspirate (FNA) biopsy samples. Features characterize the geometry of individual cell nuclei from image-derived measurements.\\

Ten nucleus properties are measured per image, each recorded as mean, standard error (SE), and worst (largest observed value across the image), yielding 30 features in total (Table~\ref{tab:feature_groups}). The target variable is binary: 0 = malignant, 1 = benign. No missing values were present in any of the 30 features.

\begin{table}[H]
\centering
\small
\setlength{\tabcolsep}{3pt}
\begin{tabular}{llp{3.8cm}}
\toprule
\textbf{Measurement} & \textbf{Description} \\
\midrule
radius    & Mean distance, center to perimeter \\
texture & Std dev of grey-scale values \\
perimeter  & Tumour perimeter \\
area       & Tumour cross-sectional area \\
smoothness & Local variation in radius lengths \\
compactness & Perimeter$^2$ / area $-$ 1 \\
concavity  & Severity of concave contour portions \\
concave pts  & Number of concave contour portions \\
symmetry   & Symmetry of the cell nucleus \\
fractal dim. & Coastline approximation $-$ 1 \\
\bottomrule
\end{tabular}
\caption{The ten nucleus measurements. Each is recorded as mean, SE, and worst, producing 30 features in total.}
\label{tab:feature_groups}
\end{table}


\subsection{Exploratory Data Analysis}

\subsubsection{Class Balance}

The dataset contains 212 malignant (37.3\%) and 357 benign (62.7\%) samples, giving a benign-to-malignant ratio of approximately 1.68\,:\,1 (Figure~\ref{fig:class_balance}). This moderate imbalance has two methodological implications. First, stratified sampling is required to preserve class proportions across train and test sets. Second, accuracy alone might be an insufficient evaluation metric, since a naive classifier that always predicts benign would achieve 62.7\% without learning any pattern.

\begin{figure}[H]
    \centering
    \includegraphics[width=0.8\linewidth]{sample_count.png}
    \caption{Class distribution of the Breast Cancer Wisconsin dataset. Benign samples (62.7\%, $n = 357$) outnumber malignant samples (37.3\%, $n = 212$), producing a moderate imbalance that motivates stratified splitting and macro-averaged evaluation metrics.}
    \label{fig:class_balance}
\end{figure}


\subsubsection{Feature Distributions by Class}

Figure~\ref{fig:distributions_all} shows class-separated distributions for four representative features per aggregate.\\

As for the mean measurements (Figure~\ref{fig:distributions_mean}), the \texttt{mean radius} shows the strongest class separation, with malignant tumors substantially larger and more variable. However, the worst measurements (Figure~\ref{fig:distributions_worst}) achieve the strongest separation. This is expected from a clinical perspective: abnormal nuclei usually indicate malignancies \citep{wolberg1995breast}, and the worst statistics evaluate this quantitatively. \texttt{worst concave points} is particularly discriminative in this aggregate.\\

On the other hand, the standard error measurements (Figure~\ref{fig:distributions_se}) are less discriminative individually. \\

Summary statistics show significant scale heterogeneity: \texttt{mean area} spans 143--2{,}501 mm$^2$, while \texttt{mean fractal dimension} is confined to 0.050--0.097. This four-order-of-magnitude range makes feature standardization essential before any algorithm is applied.

\begin{figure*}[p]
    \centering
    \begin{subfigure}{\textwidth}
        \centering
        \includegraphics[width=\textwidth]{mean.png}
        \caption{Mean measurements: \texttt{mean radius} (size), \texttt{mean texture}, \texttt{mean concavity} (shape), \texttt{mean fractal dimension} (boundary complexity).}
        \label{fig:distributions_mean}
    \end{subfigure}

    \vspace{4pt}

    \begin{subfigure}{\textwidth}
        \centering
        \includegraphics[width=\textwidth]{std.png}
        \caption{Standard error (SE) measurements: \texttt{radius error}, \texttt{texture error}, \texttt{concavity error}, \texttt{symmetry error}.}
        \label{fig:distributions_se}
    \end{subfigure}

    \vspace{4pt}

    \begin{subfigure}{\textwidth}
        \centering
        \includegraphics[width=\textwidth]{worst.png}
        \caption{Worst (largest observed) measurements: \texttt{worst radius}, \texttt{worst texture}, \texttt{worst concave points}, \texttt{worst fractal dimension}.}
        \label{fig:distributions_worst}
    \end{subfigure}

    \caption{Class-separated distributions (box plots with jittered strip) for four representative features per aggregate. Red: malignant; blue: benign. Vertical box separation indicates discriminative power.}
    \label{fig:distributions_all}
\end{figure*}



\subsubsection{Correlation Analysis}


The Pearson correlation matrix (see Figure in the Jupyter Notebook) revealed 21 feature pairs with $|r| \geq 0.90$, organized around two dominant clusters.

One comprises the geometric features (radius, perimeter, and area): the strongest pair is \texttt{mean radius} $\leftrightarrow$ \texttt{mean perimeter} ($r = 0.998$), and all nine combinations of \{radius, perimeter, area\} $\times$ \{mean, worst\} exceed $|r| = 0.94$, which is expected due to their geometric relationship.

The other cluster is related to concavity features: \texttt{mean concavity} $\leftrightarrow$ \texttt{mean concave points} ($r = 0.92$) and \texttt{mean concave points} $\leftrightarrow$ \texttt{worst concave points} ($r = 0.91$).



\subsection{Preprocessing and Feature Analysis}

\subsubsection{Stratified Train-Test Split}

The dataset is partitioned into 80\% training ($n = 455$) and 20\% test ($n = 114$) sets using stratified sampling (\texttt{stratify=y}, \texttt{random\_state=42}). Table~\ref{tab:split_balance} confirms that the class proportions are preserved.

\begin{table}[H]
\centering
\small
\setlength{\tabcolsep}{4pt}
\begin{tabular}{lccccc}
\toprule
\textbf{Split} & \textbf{n} & \textbf{Mal.} & \textbf{Ben.} & \textbf{Mal.\%} & \textbf{Ben.\%} \\
\midrule
Full       & 569 & 212 & 357 & 37.3\% & 62.7\% \\
Training   & 455 & 170 & 285 & 37.4\% & 62.6\% \\
Test       & 114 &  42 &  72 & 36.8\% & 63.2\% \\
\bottomrule
\end{tabular}
\caption{Class balance after stratified 80/20 split.}
\label{tab:split_balance}
\end{table}

The test set is held out entirely throughout all subsequent modelling steps.

\subsubsection{Hyperparameter Tuning Strategy}

All tuning is performed exclusively on the training set ($n = 455$) using stratified 5-fold cross-validation (CV). The mean CV accuracy across the five folds is used as the selection criterion. Once the optimal hyperparameters are identified, the model is on the held-out test set once.

\subsubsection{Feature Scaling}

\texttt{StandardScaler} is applied to all 30 features, transforming each to zero mean and unit variance using training-set statistics:
\[
    z = \frac{x - \hat{\mu}_{\text{train}}}{\hat{\sigma}_{\text{train}}}
\]
The scaler is \textit{fitted exclusively on the training set} and then applied identically to the test set. After scaling, the training set achieves a maximum absolute mean of $\approx 0$ and a uniform standard deviation of $\approx 1$, as expected.\\


\subsubsection{Feature Selection}
When deciding which features to keep, two criteria were used.\\
First, the variance was studied to identify variables with near-zero variance. Five features have raw variance $< 10^{-4}$: \texttt{fractal dimension error}, \texttt{smoothness error}, \texttt{concave points error}, \texttt{mean fractal dimension}, and \texttt{symmetry error}. However, since these features have inherently small absolute values and their coefficients of variation (std / $|$mean$|$) range from 11--73\%, these values indicate meaningful variability. \\

High correlation is also evaluated, as it is the most significant potential issue in this dataset and could lead to multicollinearity problems.


Nevertheless, it was decided to keep all 30 features.

Since, in this case, high intra-group correlation is an expected consequence of the measurement protocol rather than a data artifact (given the three statistical summaries per measurement), removing one member of a correlated group would discard a potentially complementary discriminative signal.


% ─────────────────────────────────────────────────────────────────────────────
\section{Classification Methods}
\label{sec:methods_class}

Four supervised classification algorithms are compared in this study: K-Nearest Neighbours (KNN), Logistic Regression, Support Vector Machines (SVM), and a Multi-Layer Perceptron (Neural Network). Table~\ref{tab:model_comparison} summarises their key assumptions, strengths, and limitations as reported in the literature \citep{uddin2019comparing, hastie2009elements, bishop2006pattern, james2021statistical}.

\begin{table*}[t]
\centering
\small
\setlength{\tabcolsep}{5pt}
\renewcommand{\arraystretch}{1.25}
\begin{tabularx}{\textwidth}{p{0.15\textwidth}XXX}
\toprule
\textbf{Algorithm} & \textbf{Assumptions} & \textbf{Strengths} & \textbf{Limitations} \\
\midrule
\textbf{K-Nearest Neighbors}
  & Instances close in feature space belong to the same class; all features contribute equally; distance is meaningful (requires scaling).
  & Simple to implement; no explicit training phase; naturally handles multi-class problems; can model non-linear boundaries.
  & Prediction requires computing distances to all training points, making inference cheap at this dataset size ($n=455$) but increasingly costly as $n$ grows. Sensitive to irrelevant or correlated features; provides no information on feature importance. \\
\addlinespace
\textbf{Logistic Regression}
  & Linear relationship between features and the log-odds of the outcome; observations are independent; no severe multicollinearity.
  & Highly interpretable coefficients; probabilistic outputs; fast to train and update; no distributional assumption on predictors.
  & Cannot capture complex non-linear feature interactions; tends to underfit if the true boundary is non-linear; vulnerable to sampling bias. \\
\addlinespace
\textbf{SVM (RBF)}
  & Optimal decision boundary lies in a transformed (kernel) feature space; data are separable with a margin (soft-margin via $C$).
  & Strong generalization through margin maximization; robust to high-dimensional spaces; kernel trick allows non-linear boundaries without explicit feature maps; less prone to overfitting than many methods.
  & Computationally expensive for large datasets; sensitive to kernel and hyperparameter choices; the RBF kernel maps data into an implicit high-dimensional space, making the decision boundary difficult to interpret in terms of the original features. \\
\addlinespace
\textbf{Neural Network (MLP)}
  & No explicit distributional assumptions; sufficient data and network capacity can approximate any continuous function (universal approximation).
  & Can detect complex non-linear relationships and feature interactions; flexible architecture; multiple training algorithms available.
  & Black-box model with limited interpretability; computationally expensive to train; sensitive to hyperparameter choices (architecture, learning rate, regularisation); requires careful preprocessing.\\
\bottomrule
\end{tabularx}
\caption{Comparison of the four supervised classification algorithms evaluated in this study. Strengths and limitations are consistent with findings from a broad review of supervised machine learning for disease prediction \citep{uddin2019comparing}.}
\label{tab:model_comparison}
\end{table*}

\subsection{K-Nearest Neighbours}
\label{sec:knn}

K-Nearest Neighbours (KNN) is a classification algorithm that assigns a
class to a data point based on the majority vote of its $k$ nearest
neighbours in the dataset, using Euclidean distance by default
\citep{hastie2009elements}. Since KNN relies on distance calculations,
it is sensitive to feature scale, motivating the \texttt{StandardScaler}
applied in preprocessing.

The choice of $k$, the only hyperparameter of the model, has a significant
effect on the classifier obtained \citep{james2021statistical}: small
values of $k$ produce an overly flexible decision boundary that captures
noise rather than the true underlying pattern, corresponding to low bias
but very high variance. As $k$ grows, the boundary becomes less flexible
and closer to linear, corresponding to low variance but high bias. The
optimal $k$ was therefore determined empirically by sweeping
$k \in \{1, \ldots, 20\}$ and recording the mean 5-fold CV accuracy on the
training set for each value (Figure~\ref{fig:knn_sweep}).

\begin{figure}[H]
    \centering
    \includegraphics[width=\linewidth]{KNN.png}
    \caption{Mean 5-fold CV accuracy on the training set as a function of $k$ ($\pm1$ std shaded).
$k = 9$ achieves the highest CV accuracy among odd candidates (0.9692)
with the lowest standard deviation (0.0162), and is selected as optimal.}
    \label{fig:knn_sweep}
\end{figure}

Among odd values of $k$ (required to avoid ties in binary majority voting), $k = 3$, $k = 7$, and $k = 9$ all achieve the maximum mean CV accuracy of 96.92\%. $k = 9$ was selected as the optimal value: it matches the highest CV accuracy among odd candidates and yields the lowest standard deviation (0.0162 vs 0.0245 for $k=3$ and 0.0213 for $k=7$), indicating greater stability across folds. This is particularly relevant in a clinical context, where consistent performance across data subsets is preferable to a single high-performing split.

The sweep also reveals a clear bias-variance pattern:
$k = 1$ exhibits the highest variance (overfitting), performance stabilises between
$k = 3$ and $k = 9$, and gradually declines for larger $k$ as the boundary becomes
overly smooth (underfitting).


\subsection{Logistic Regression}
\label{sec:lr}

Logistic Regression is a linear probabilistic classification model that estimates the probability of class membership through the logistic function. In this study, it serves as an interpretable baseline for breast cancer classification, providing a transparent reference against which more flexible models can be compared.

Each coefficient indicates how a standardized feature contributes to the predicted probability of malignancy or benignity. This is especially relevant for the Wisconsin Diagnostic Breast Cancer dataset, where recent work has shown that Logistic Regression remains highly effective at distinguishing benign from malignant tumours despite its relative simplicity \citep{cheng2025lr}.

In clinical decision-support settings, this balance between accuracy and explainability is particularly valuable. The model's probabilistic outputs can support risk-based interpretation, while its transparent coefficients help identify the tumour characteristics that most strongly influence the classification.

\subsubsection{Hyperparameter tuning ($C$ parameter sweep)}

The regularisation parameter $C$ was tuned by performing a grid search over logarithmically spaced values in the range:
\[
C \in [10^{-3}, 10^{3}]
\]

Smaller values of $C$ correspond to stronger regularisation, while larger values allow more complex models.
The model was evaluated using 5-fold CV on the training set for each value of $C$.


\begin{figure}[H]
\centering
\includegraphics[width=\linewidth]{LR.png}
\caption{Mean 5-fold CV accuracy versus regularisation parameter $C$ (log scale). The best mean CV performance is achieved at $C = 1$.}
\label{fig:lr_accuracy_c}
\end{figure}

Following the principle of parsimony, the smallest value achieving the maximum mean CV accuracy was selected: \textbf{$C = 1$}.



\subsection{Support Vector Machines}
\label{sec:svm}

SVMs find the maximum-margin hyperplane separating the two classes; the
\emph{kernel trick} allows non-linear boundaries without explicitly computing
the feature map \citep{hastie2009elements}. Two key hyperparameters govern the
RBF kernel: regularisation strength $C$ and bandwidth $\gamma$
\citep{james2021statistical}. All cross-validation uses \texttt{balanced\_accuracy}
as the scoring metric, giving equal weight to both classes.

\subsubsection{Kernel Comparison}

Linear, RBF, and Polynomial ($C = 1$, degree $= 3$) kernels were compared using 5-fold CV on the training set at default settings to obtain a baseline ranking (Table~\ref{tab:svm_kernels}). The Polynomial kernel underperformed markedly
(86.18\% CV balanced accuracy), suggesting the decision boundary is not well modelled by a
cubic surface. RBF achieved 96.89\% and Linear achieved 96.30\% CV balanced accuracy; RBF was selected for
tuning due to its greater flexibility.

\begin{table}[H]
\centering
\small
\setlength{\tabcolsep}{3pt}
\begin{tabular}{lccc}
\toprule
\textbf{Kernel} & \textbf{CV Bal. Acc} & \textbf{CV AUC} & \textbf{Time (s)} \\
\midrule
Linear          & 0.9630 & 0.9916 & 0.004 \\
RBF             & 0.9689 & 0.9938 & 0.005 \\
Polynomial      & 0.8618 & 0.9914 & 0.004 \\
\bottomrule
\end{tabular}
\caption{Mean 5-fold CV balanced accuracy and CV AUC on the training set for each kernel at default $C = 1$. RBF and Linear kernels perform similarly; Polynomial underperforms markedly.}
\label{tab:svm_kernels}
\end{table}

\subsubsection{RBF Hyperparameter Search}

\texttt{GridSearchCV} with 5-fold CV was run on the training set over
$C \in \{0.01,\,0.1,\,1,\,10,\,100\}$ and
$\gamma \in \{0.001,\,0.01,\,0.1,\,1,\,10\}$ (25 combinations, 125 fits).
The optimal pair is $C = 10$, $\gamma = 0.01$ (mean CV balanced accuracy 97.59\%).
CV balanced accuracy collapses for $\gamma \geq 1$ due to over-localised support vectors
\citep{hastie2009elements} (Figures~\ref{fig:svm_grid} and \ref{fig:svm_c_gamma}).

\begin{figure}[H]
    \centering
    \includegraphics[width=0.8\linewidth]{SVM_quadre_1.png}
    \caption{Mean 5-fold CV balanced accuracy for all 25 $(C, \gamma)$ combinations. Cyan border: best cell ($C = 10$, $\gamma = 0.01$).}
    \label{fig:svm_grid}
\end{figure}

\begin{figure}[H]
    \centering
    \includegraphics[width=0.8\linewidth]{In_curv_2.png}
    \includegraphics[width=0.8\linewidth]{In_curv_1.png}
    \caption{CV balanced accuracy vs $\gamma$ (top, $C = 10$) and vs $C$
    (bottom, $\gamma = 0.01$). The optimal region is $\gamma \leq 0.01$.}
    \label{fig:svm_c_gamma}
\end{figure}


\subsection{Neural Network}
\label{sec:nn}

A Multi-Layer Perceptron (MLP) is a feedforward neural network that learns non-linear decision boundaries through a series of fully connected layers, each applying a weighted linear transformation followed by a non-linear activation function \citep{bishop2006pattern}. Unlike the previous methods, MLPs do not assume a fixed functional form for the decision boundary, making them capable of capturing complex feature interactions at the cost of reduced interpretability and increased sensitivity to hyperparameter choices \citep{hastie2009elements}.

\subsubsection{Architecture Selection}

Two hidden-layer configurations were evaluated on the preprocessed dataset (455 training samples, 30 standardised features):

\begin{itemize}
    \item \textbf{Architecture 1: (64,):} one hidden layer with 64 neurons, selected as a shallow baseline larger than the 30 input features but modest enough to avoid overfitting on this dataset size.
    \item \textbf{Architecture 2: (64, 32):} two hidden layers in a funnel structure. The first layer learns broad feature combinations from the 30 inputs; the second compresses these into higher-level representations before the output, reflecting the clinical intuition that malignancy is determined by combinations of measurements rather than any single feature.
\end{itemize}

ReLU activation was used in all hidden layers to avoid the vanishing gradient problem associated with sigmoid activations \citep{hastie2009elements}. The Adam optimiser was selected with its default learning rate ($\eta = 0.001$). Two regularisation strategies were applied: L2 weight decay via the sklearn default (alpha $= 0.0001$), and early stopping with a validation fraction of 10\% and patience of 20 epochs (\texttt{n\_iter\_no\_change} $= 20$), which halts training automatically when validation performance stops improving.

\subsubsection{Architecture Comparison and Loss Curves}

\begin{table}[H]
\centering
\small
\setlength{\tabcolsep}{4pt}
\begin{tabular}{lcc}
\toprule
\textbf{Architecture} & \textbf{CV Accuracy} & \textbf{Training Time} \\
\midrule
(64,) 1 hidden layer     & 0.9296 & $\sim$0.11s \\
(64, 32) 2 hidden layers & 0.9465 & $\sim$0.42s \\
\bottomrule
\end{tabular}
\caption{Architecture comparison based on mean 5-fold CV balanced accuracy on the training set. The (64, 32) architecture achieves higher balanced accuracy and lower variance across folds.}
\label{tab:mlp_arch}
\end{table}

Architecture (64, 32) outperformed the single hidden layer by 1.7 percentage points in mean CV balanced accuracy (0.9465 vs 0.9296), with lower variance across folds ($\pm$0.0244 vs $\pm$0.0311). Figure~\ref{fig:mlp_loss} shows the training loss and validation score per epoch for both architectures. Both converged smoothly within 42 to 52 epochs, with training loss decreasing steadily and validation score plateauing near 1.0, indicating healthy generalisation with no overfitting.

\begin{figure}[H]
    \centering
    \includegraphics[width=0.8\linewidth]{NN_1.png}
    \includegraphics[width=0.8\linewidth]{NN_2.png}
    \caption{Training loss and validation score per epoch for both MLP architectures. Both converge smoothly with no overfitting.}
    \label{fig:mlp_loss}
\end{figure}

\subsubsection{L2 Regularisation Verification}

A sweep over alpha $\in \{0.0001, 0.01, 1.0, 10.0\}$ was applied to the best architecture to confirm the default sklearn value is appropriate. Training balanced accuracy was monitored across values: alpha $= 0.0001$ and $0.01$ both maintained high balanced accuracy (0.97 to 0.98), while alpha $= 1.0$ began to over-constrain the model and alpha $= 10.0$ caused underfitting (balanced accuracy $\sim$0.85). The default alpha $= 0.0001$ was confirmed as optimal.


% ─────────────────────────────────────────────────────────────────────────────
\section{Results}
\label{sec:results}

The final evaluation results for all four classifiers on the held-out test set ($n = 114$) are evaluated. \footnote{\textit{Note on terminology:} Throughout this notebook, clinical language is used consistently: a false negative refers to a malignant case predicted as benign (missed cancer diagnosis), and a false positive refers to a benign case predicted as malignant (unnecessary alarm). This follows clinical convention where positive = cancer present, regardless of the dataset encoding (malignant=0, benign=1). ROC curves are computed with benign as the positive class (sklearn default given the encoding) and should be interpreted accordingly. 
}

\subsection{KNN: Final Evaluation}

The final KNN model ($k = 9$) correctly classified 111 of 114 test samples,
achieving a test accuracy of 97.37\%. All 72 benign cases were identified
correctly (recall\,=\,1.000), while 3 of 42 malignant cases were misclassified
as benign (recall\,=\,0.929). These 3 false negatives likely correspond to borderline cases lying close
to the decision boundary in feature space; in a clinical screening context,
they carry a higher cost than false positives as they represent missed diagnoses.
Table~\ref{tab:knn_report} summarises the full classification report.




\begin{table}[H]
\centering
\small
\setlength{\tabcolsep}{4pt}
\begin{tabular}{lcccc}
\toprule
\textbf{Class} & \textbf{Precision} & \textbf{Recall} & \textbf{F1} & \textbf{Support} \\
\midrule
Malignant    & 1.000 & 0.929 & 0.963 & 42 \\
Benign       & 0.960 & 1.000 & 0.980 & 72 \\
\midrule
Accuracy     & &  & {0.974} & 114 \\
Macro avg    & 0.980 & 0.964 & 0.971 & 114 \\
Weighted avg & 0.975 & 0.974 & 0.974 & 114 \\
\bottomrule
\end{tabular}
\caption{Classification report for the final KNN model ($k = 9$).
Macro F1\,=\,0.971 reflects strong performance on both classes,
with zero false positives and three false negatives.}
\label{tab:knn_report}
\end{table}



\subsection{Logistic Regression: Final Evaluation}

The final Logistic Regression model ($C = 1$) achieved an accuracy of 98.25\%, correctly distinguishing most malignant and benign cases. Precision, recall, and F1-score were highly balanced across both classes, indicating that the classifier did not favour the majority benign class despite the moderate class imbalance.

\begin{table}[H]
\centering
\small
\setlength{\tabcolsep}{4pt}
\begin{tabular}{lcccc}
\toprule
\textbf{Class} & \textbf{Precision} & \textbf{Recall} & \textbf{F1} & \textbf{Support} \\
\midrule
Malignant    & 0.976 & 0.976 & 0.976 & 42 \\
Benign       & 0.986 & 0.986 & 0.986 & 72 \\
\midrule
Accuracy     & & & {0.983} & 114 \\
Macro avg    & 0.981 & 0.981 & 0.981 & 114 \\
Weighted avg & 0.983 & 0.983 & 0.983 & 114 \\
\bottomrule
\end{tabular}
\caption{Final test-set classification report for Logistic Regression ($C = 1$). The model achieves 98.25\% accuracy with similarly high precision, recall, and F1-score for both classes.}
\label{tab:lr_report}
\end{table}

\subsection{SVM: Final Evaluation}

The final RBF SVM ($C = 10$, $\gamma = 0.01$) correctly classified 112 of 114 samples, with one false negative and one false positive, giving a more balanced error profile than KNN. Table~\ref{tab:svm_report} reports the full metrics; AUC = 0.9977.

\begin{table}[H]
\centering
\small
\setlength{\tabcolsep}{4pt}
\begin{tabular}{lcccc}
\toprule
\textbf{Class} & \textbf{Precision} & \textbf{Recall} & \textbf{F1} & \textbf{Support} \\
\midrule
Malignant    & 0.9762 & 0.9762 & 0.9762 & 42 \\
Benign       & 0.9861 & 0.9861 & 0.9861 & 72 \\
\midrule
Accuracy     & &  & {0.9825} & 114 \\
Macro avg    & 0.9812 & 0.9812 & 0.9812 & 114 \\
Weighted avg & 0.9825 & 0.9825 & 0.9825 & 114 \\
\bottomrule
\end{tabular}
\caption{Classification report for the RBF SVM ($C=10$, $\gamma=0.01$,
hyperparameters from 5-fold CV balanced accuracy). AUC = 0.9977.}
\label{tab:svm_report}
\end{table}


Feature importance was assessed from the standardized Logistic Regression coefficients. Because the target variable is encoded as 0 for malignant tumours and 1 for benign tumours, negative coefficients increase the probability of the malignant class, whereas positive coefficients are associated with benign predictions. Therefore, the features with the largest negative coefficients can be interpreted as the strongest indicators of malignancy in this model.

The most influential malignancy-associated features were mainly worst-case measurements, suggesting that the largest observed nuclear abnormalities carry substantial diagnostic information. The top features were:

\begin{itemize}[noitemsep, topsep=0pt, label=\tiny$\bullet$]
    \item \texttt{worst texture}
    \item \texttt{radius error}
    \item \texttt{worst area}
    \item \texttt{area error}
    \item \texttt{worst radius}
    \item \texttt{worst concave points}
    \item \texttt{worst symmetry}
    \item \texttt{worst concavity}

\end{itemize}

Overall, this pattern is consistent with the exploratory analysis: malignant tumours tend to present larger, more irregular, and more heterogeneous nuclei. The predominance of worst-case features also supports the clinical relevance of extreme nuclear measurements when distinguishing malignant from benign samples.


\subsection{Neural Network: Final Evaluation}

The final MLP (64, 32) achieved 97.37\% test accuracy and AUC = 0.996. Malignant recall of 0.95 indicates that 2 of 42 malignant cases were misclassified as benign (false negatives). In a clinical screening context, false negatives represent missed diagnoses and are more costly than false positives. Table~\ref{tab:mlp_report} summarises the metrics.

\begin{table}[H]
\centering
\small
\setlength{\tabcolsep}{4pt}
\begin{tabular}{lcccc}
\toprule
\textbf{Class} & \textbf{Precision} & \textbf{Recall} & \textbf{F1} & \textbf{Support} \\
\midrule
Malignant    & 0.976 & 0.952 & 0.964 & 42 \\
Benign       & 0.973 & 0.986 & 0.979 & 72 \\
\midrule
Accuracy     & &  & {0.974} & 114 \\
Macro avg    & 0.974 & 0.969 & 0.972 & 114 \\
Weighted avg & 0.974 & 0.974 & 0.974 & 114 \\
\bottomrule
\end{tabular}
\caption{Classification report for the final MLP model (64, 32). Macro F1 = 0.97, AUC = 0.996.}
\label{tab:mlp_report}
\end{table}

\subsection{Confusion Matrix Comparison}

Figure~\ref{fig:cm_all} shows the confusion matrices for all four models, allowing direct comparison of error patterns. KNN produces zero false positives but three false negatives. Logistic Regression and SVM each make two errors: LR has one FN and one FP; SVM has one FN and one FP. The MLP produces two false negatives and one false positive.

\begin{figure*}[t]
    \centering
    \begin{subfigure}{0.46\textwidth}
        \centering
        \includegraphics[width=0.8\linewidth]{CM_KNN.png}
        \caption{KNN ($k=9$)}
        \label{fig:cm_knn}
    \end{subfigure}\hfill
    \begin{subfigure}{0.46\textwidth}
        \centering
        \includegraphics[width=0.8\linewidth]{CM_LR.png}
        \caption{Logistic Regression ($C=1$)}
        \label{fig:cm_lr}
    \end{subfigure}

    \vspace{8pt}

    \begin{subfigure}{0.46\textwidth}
        \centering
        \includegraphics[width=0.8\linewidth]{CM_SVM.png}
        \caption{SVM (RBF, $C=10$, $\gamma=0.001$)}
        \label{fig:cm_svm}
    \end{subfigure}\hfill
    \begin{subfigure}{0.46\textwidth}
        \centering
        \includegraphics[width=0.8\linewidth]{CM_NN.png}
        \caption{MLP (64, 32)}
        \label{fig:cm_mlp}
    \end{subfigure}
    \caption{Confusion matrices for all four classifiers on the held-out test set ($n = 114$). Rows: true class; columns: predicted class. All models achieve similar accuracy; differences are visible in the distribution of false negatives (malignant cases missed) and false positives (benign cases incorrectly flagged).}
    \label{fig:cm_all}
\end{figure*}

\begin{figure}[H]
    \centering
    \includegraphics[width=\linewidth]{ROC_curves.jpeg}
    \caption{ROC curves for all four classifiers on the held-out test set.
    All models achieve AUC $\geq 0.994$, indicating near-perfect discrimination
    between malignant and benign cases across all classification thresholds.}
    \label{fig:roc_all}
\end{figure}



% ─────────────────────────────────────────────────────────────────────────────
\section{Discussion}
\label{sec:discussion}


\subsection{Best Performing Model}

Logistic Regression achieved the highest accuracy of approx 98\%, making it difficult to declare a single winner on performance alone. However, we considered Logistic Regression as the best overall model for this task since it matches SVM 
on accuracy and F1, trains faster, and offers direct coefficient interpretability that SVM cannot provide. It produces directly interpretable coefficients that clinicians can use to understand which tumour features drive the prediction \citep{bishop2006pattern}. In a clinical decision-support context, interpretability and computational efficiency are as important as raw accuracy \citep{james2021statistical}.

\subsection{Dataset Properties and Algorithm Choice}

Three properties of this dataset significantly influenced algorithm selection and performance.

\textbf{High dimensionality (30 features).} With 30 features, distance-based methods like KNN are susceptible to the curse of dimensionality, where distances become less meaningful in high-dimensional spaces. Feature scaling was mandatory to prevent features with large absolute ranges from dominating distance calculations \citep{hastie2009elements}. Linear models like Logistic Regression handle high dimensionality well through regularisation.


\textbf{Moderate class imbalance (63\% benign, 37\% malignant).} The imbalance motivated splitting to preserve class proportions across train and test sets.


\subsection{Limitations}

Despite the strong results, this analysis has several limitations.

All 30 features were retained despite high correlation among geometric features (radius, perimeter, area). While this ensures a fair comparison across models, future work could apply dimensionality reduction such as PCA to remove redundant information and improve model efficiency \citep{james2021statistical}.

The dataset comes from a single institution with only 569 samples. Performance may not generalise to patients from different hospitals or demographic backgrounds, which is a common challenge in clinical machine learning \citep{james2021statistical}.
For the Neural Network, only two architectures were compared. A broader search over layer sizes and regularisation values would provide stronger justification for the final configuration.


% ─────────────────────────────────────────────────────────────────────────────
\section{Conclusions}
\label{sec:conclusions}
All four models achieved high performance, with test accuracies between 97.37\% (KNN and Neural Network) and 98.25\% (Logistic Regression and SVM). Hyperparameters were selected using 5-fold cross-validation on the training set for all models, ensuring the test set was used only once for final evaluation.

Logistic Regression is the recommended model for clinical use. It matches the top accuracy of SVM, trains in milliseconds, and its coefficients are directly interpretable, clinicians can see which tumour features drive each prediction \citep{bishop2006pattern}. This combination of performance, speed, and transparency is difficult to match in a healthcare setting.

Across all models, malignant recall was 0.93 or above, meaning at most 7\% of cancer cases were missed. Future work should explore adjusting the classification threshold to further reduce false negatives, and test model performance on data from other institutions \citep{hastie2009elements}.

% ─────────────────────────────────────────────────────────────────────────────
\section*{AI Usage}
For this assignment, we used AI tools to help optimise and elaborate our code and improve the overall structure and flow of the text. That said, all the core decisions and analytical approaches are our own, based primarily on the class materials and literature.



%─────────────────────────────────────────────────────────────────────────────
\bibliographystyle{plainnat}
\begin{thebibliography}{99}

\bibitem[Bishop(2006)]{bishop2006pattern}
Bishop, C.~M. (2006).
\newblock \textit{Pattern Recognition and Machine Learning}.
\newblock Springer.

\bibitem[Hastie et~al.(2009)]{hastie2009elements}
Hastie, T., Tibshirani, R., \& Friedman, J. (2009).
\newblock \textit{The Elements of Statistical Learning} (2nd ed.).
\newblock Springer.

\bibitem[James et~al.(2021)]{james2021statistical}
James, G., Witten, D., Hastie, T., \& Tibshirani, R. (2021).
\newblock \textit{An Introduction to Statistical Learning} (2nd ed.).
\newblock Springer.

\bibitem[Pedregosa et~al.(2011)]{pedregosa2011sklearn}
Pedregosa, F., et~al. (2011).
\newblock Scikit-learn: Machine learning in Python.
\newblock \textit{Journal of Machine Learning Research}, 12, 2825--2830.

\bibitem[Siegel et~al.(2023)]{siegel2023cancer}
Siegel, R.~L., Miller, K.~D., Wagle, N.~S., \& Jemal, A. (2023).
\newblock Cancer statistics, 2023.
\newblock \textit{CA: A Cancer Journal for Clinicians}, 73(1), 17--48.

\bibitem[Wolberg et~al.(1995)]{wolberg1995breast}
Wolberg, W.~H., Street, W.~N., \& Mangasarian, O.~L. (1995).
\newblock Breast cytology diagnoses via digital image analysis.
\newblock \textit{Analytical and Quantitative Cytology and Histology},
17(2), 76--84.

\bibitem[Nadkarni(2016)]{nadkarni2016datamining}
Nadkarni, P. (2016).
\newblock Core technologies: Data mining and `big data'.
\newblock In \textit{Clinical Research Computing} (pp. 187--204).
\newblock Elsevier.
\newblock \url{https://doi.org/10.1016/B978-0-12-803130-8.00010-5}

\bibitem[Cheng and Yu(2025)]{cheng2025lr}
Cheng, W., \& Yu, Z. (2025).
\newblock A fast and interpretable logistic regression framework for breast tumor classification using the Wisconsin Diagnostic Dataset.
\newblock \textit{medRxiv}.
\newblock \url{https://doi.org/10.64898/2025.12.23.25342946}

\bibitem[Uddin et~al.(2019)]{uddin2019comparing}
Uddin, S., Khan, A., Hossain, M.~E., \& Moni, M.~A. (2019).
\newblock Comparing different supervised machine learning algorithms for disease prediction.
\newblock \textit{BMC Medical Informatics and Decision Making}, 19(1), 281.
\newblock \url{https://doi.org/10.1186/s12911-019-1004-8}

\end{thebibliography}

\end{document}